% =============================================================================
% Methodology
% =============================================================================
\section{Methodology}\label{sec:methodology}

\subsection{Numerical Jacobian}

We compute the Jacobian in \eqref{eq:jacobian_raw} by central finite
differences.  For parameter $\theta_j$,
\begin{equation}\label{eq:fd}
  \frac{\partial f_i}{\partial \theta_j}
  \approx
  \frac{f_i(\boldsymbol{\theta}+\delta_j\mathbf{e}_j)
       -f_i(\boldsymbol{\theta}-\delta_j\mathbf{e}_j)}
       {2\delta_j},
\end{equation}
where $\mathbf{e}_j$ is the unit vector in the $j$th parameter direction and
\begin{equation}
  \delta_j = 10^{-6}\,|\theta_j|.
\end{equation}
Step-size convergence was checked against neighbouring perturbation levels; the
reported Jacobians were insensitive to those changes at the quoted precision.

For identifiability analysis we use the scaled Jacobian
\eqref{eq:jacobian_scaled}, not the raw Jacobian.  This distinction is
essential.  Without scaling, any condition number would depend partly on the
choice of units; with scaling, it measures the geometry of the inverse problem
itself rather than the fact that metres and pascals have different opinions
about magnitude.

\subsection{Newton-type inversion}

Given synthetic or measured frequencies $\mathbf{f}^\mathrm{obs}$, we recover
$\boldsymbol{\theta}=(a,c,E)^\top$ by solving a nonlinear least-squares
problem.  The residual vector is normalised by the observed frequencies,
\begin{equation}\label{eq:residual}
  r_i(\boldsymbol{\theta})
  =
  \frac{f_i(\boldsymbol{\theta})-f_i^\mathrm{obs}}
       {f_i^\mathrm{obs}},
\end{equation}
and we minimise
\begin{equation}\label{eq:cost}
  \Phi(\boldsymbol{\theta})
  =
  \frac{1}{2}\sum_i r_i^2.
\end{equation}
The solver is a Newton-type trust-region least-squares iteration applied in
normalised parameter coordinates,
\begin{equation}
  x_j = \frac{\theta_j}{\theta_{j,0}},
\end{equation}
where $\theta_{j,0}$ is the initial guess.  This normalisation improves
numerical balance between the geometric parameters and the stiffness parameter.

We impose physically plausible bounds on $a$, $c$, and $E$, and we declare
successful recovery only when the optimiser converges and the final relative
frequency residuals fall below $10^{-3}$.  Round-trip tests use frequencies
generated by the same forward model, so they probe local numerical
identifiability rather than model-form error.

\subsection{Three-mode and five-mode inversion}

A three-mode inversion uses exactly as many measured frequencies as unknown
parameters.  This is sufficient for a local solve, but it is fragile.  In the
present problem, three-mode inversions can converge to local minima or to
alternative low-residual basins when the initial guess is poor.  We therefore
distinguish between:
\begin{itemize}
  \item \textbf{local recovery}, assessed by round-trip inversion near the true
  parameter set; and
  \item \textbf{global robustness}, assessed by repeated solves from perturbed
  initial guesses.
\end{itemize}
In practice, using at least five modes ($n=2,\ldots,6$) makes the problem
overdetermined and removes the observed failure modes.  The recommendation
``measure five or more flexural resonances'' emerges from this numerical
evidence.

\subsection{Condition-number maps}

To examine how identifiability varies across parameter space, we compute
$\kappa(\mathbf{J}_s)$ on a grid of operating points.  It is convenient to
replace $c$ by a geometric flattening parameter
\begin{equation}
  \zeta_g = 1 - \frac{c}{a},
\end{equation}
where the subscript $g$ distinguishes geometric flattening from the damping
ratio.  We then examine representative slices of the map in the
$(a,\zeta_g)$ plane while varying $E$ over the interval
\SIrange{0.05}{0.5}{MPa}.  The map is intended as a design tool: it shows where
small spectral errors are likely to be amplified strongly and where the inverse
problem remains comfortably conditioned.

\subsection{Noise model and uncertainty quantification}

For the Cram\'er--Rao analysis we assume independent Gaussian frequency errors
with common fractional standard deviation $\epsilon_\mathrm{noise}$.  The
results are reported mainly for $\epsilon_\mathrm{noise}=0.01$, corresponding to
$1\%$ frequency uncertainty.  Because the Fisher matrix in
\eqref{eq:fisher} scales as $\epsilon_\mathrm{noise}^{-2}$, the corresponding
standard deviations scale linearly with noise level.  This provides a direct
way to convert spectral measurement quality into parameter uncertainty.

\subsection{Cross-application case}

To test whether the conclusions are specific to the canonical abdomen, we repeat
the analysis for a ripe watermelon, parameterised by
$a=\SI{0.158}{m}$, $c=\SI{0.123}{m}$, $h=\SI{0.015}{m}$, and
$E=\SI{50}{MPa}$, with the remaining properties taken from the corresponding
material model in the repository.  This case is useful because it preserves the
same inverse structure while moving to a very different stiffness scale.

\subsection{Prolate comparison}

For the prolate comparison, we parameterise the shell with semi-minor axis
$a < $ semi-major axis $c$, sweeping eccentricity
$\varepsilon = \sqrt{1 - (a/c)^2}$ over $[0.3, 0.65]$ at fixed
equivalent-sphere radius $R_\mathrm{eq} = (a^2 c)^{1/3}$.  The same five-mode
Ritz framework ($n=2,\ldots,6$) and Jacobian analysis are applied.  Because the
curvature distribution is inverted relative to the oblate case --- maximal at the
poles rather than the equator --- the prolate sweep tests whether identifiability
restoration is a universal consequence of symmetry breaking or a property
specific to the oblate curvature--mode relationship.
